\documentclass{revtex4}
\usepackage[utf8]{inputenc}
\usepackage[T1]{fontenc}
\usepackage{geometry}
\usepackage{graphicx}
\usepackage{float}
\usepackage{hyperref}
\usepackage[english,indonesian]{babel}

\title{Optimasi Pemrosesan Template LaTeX Menggunakan Pendekatan Berbasis AI}
\author{Academic Support Agent}
\date{\today}

\begin{document}
\maketitle

\begin{abstract}

\end{abstract}

\chapter{PENDAHULUAN}


\chapter{TINJAUAN PUSTAKA}


% Original content (commented out):
% %%%%%%%%%%%%%%%%%%%%%%%%%%%%%%%%%%%%%%%%%%%%%%%%%%%%%%%%%%%%%%%%%%%%%%%%%%%
%%%%%%%%%%%%%%%%%%%%%%%%%%%%%%%%%%%%%%%%%%%%%%%%%%%%%%%%%%%%%%%%%%%%%%%%%%%
%%%%                                                                   %%%%
%%%%  MODELO LaTeX PARA UM RELATÓRIO IYPT-BRASIL                       %%%%
%%%%                                                                   %%%%
%%%%  AUTOR:  Fabricio Marques do Carmo, fabriciomarques@if.usp.br     %%%%
%%%%                                                                   %%%%
%%%%  VERSÃO: Beta        DATA: 13/08/2023        Use com sabedoria!   %%%%
%%%%                                                                   %%%%
%%%%%%%%%%%%%%%%%%%%%%%%%%%%%%%%%%%%%%%%%%%%%%%%%%%%%%%%%%%%%%%%%%%%%%%%%%%
%%%%%%%%%%%%%%%%%%%%%%%%%%%%%%%%%%%%%%%%%%%%%%%%%%%%%%%%%%%%%%%%%%%%%%%%%%%

% ATENÇÃO!!!  %%%%%%%%%%%%%%%%%%%%%%%%%%%%%%%%%%%%%%%%%%%%%%%%%%%%%%%%%%%%%
% Para melhor funcionamento no Overleaf, use: 
% compiler: XeLaTeX / TeXLive version: 2019 

%  PREÂMBULO  %%%%%%%%%%%%%%%%%%%%%%%%%%%%%%%%%%%%%%%%%%%%%%%%%%%%%%%%%%%%%  

% Formatação geral  ======================================================
\documentclass[a4paper,brazil,12pt,nofootinbib,aps,pra]{revtex4}

% Idioma e escrita========================================================
\usepackage[brazil]{babel}   % brasileiros usam acentos, cedilhas, etc.
\usepackage[T1]{fontenc}
% \usepackage[latin1]{inputenc}
\usepackage{moreverb}        % configuração de verbatim

% Símbolos matemáticos e equações  =======================================
\usepackage{amsmath}  % Fonte de alguns caracteres matemáticos
\usepackage{slashed}
\def\slash#1{#1\!\!\! /}
\usepackage{bm}       % Bold math
\usepackage{braket}   % \bra{}, \ket{} e \braket{} para Física Quântica
\usepackage{nicefrac} % \nicefrac{}{} para frações estilizadas em declive
\usepackage{bbm}      % Fonte \mathbbm{} do símbolo 1 da matriz identidade 
                      % e do R estilizado do conjunto dos números reais.
\usepackage{mathrsfs} % Fonte \mathscr{} do símbolo L da densidade lagrangiana.
\usepackage{multirow} % Células com múltiplas linhas ou colunas em tabelas

% Formatação geral  ======================================================
\usepackage{xcolor}

% Configurações de operadores matemáticos  ===============================
\renewcommand{\sin}{\operatorname{sen}}   % escreve o símbolo de seno em 
                                          % português
\renewcommand{\csc}{\operatorname{cosec}} % escreve o símbolo de cosecante em 
                                          % português
%\renewcommand{\tan}{\operatorname{tg}}    % escreve o símbolo de tangente em 
                                          % português
\DeclareMathOperator{\tr}{Tr}             % define o operador traço
\DeclareMathOperator{\diag}{diag}         % define o operador diagonal principal

\newcommand{\imineq}[2]{\vcenter{\hbox{\includegraphics[height=#2ex]{#1}}}} 
                                          % imagens em equações, pois quem não 
                                          % curte uns diagramas de Feynman?
                                          
% Formato das figuras  ==================================================  
\usepackage{graphicx}
\usepackage{subfigure}
\newcommand{\incps}[5]{\includegraphics[#2,#3][#4,#5]{#1}}
\newcommand{\incpicwh}[3]{\includegraphics[width=#2,height=#3]{#1}}
\usepackage{caption}
\DeclareCaptionFont{xipt}{\fontsize{11}{13}\mdseries}
\usepackage[font=xipt,labelfont=bf]{caption}

% Formato das tabelas  ==================================================
%\usepackage{dcolumn} %alinhar colunas das tabelas em pontos decimais


%  Formato de numeração das equações  ===================================
%\makeatletter
%\@addtoreset{equation}{subsection}
%\renewcommand{\theequation}{\thesection.\thesubsection.\arabic{equation}}
%\makeatother

% Especificações do regulamento IYPT-Brasil  ============================
\usepackage[scaled]{helvet}  % usar fonte Arial (Helvetica)              
\renewcommand*\familydefault{\sfdefault}  
\usepackage{setspace}        % controle do espaçamento entre linhas
                             % depois deve ser usado o comando \setstretch{2}
                             % dentro do texto para espaçamento duplo

% Hifenização: corrige erros de hifenização nas quebras de linha  =======
\hyphenation{tem-pe-ra-tu-ra}

% Bold symbol macro for standard LaTeX users  ===========================
\makeatletter
\renewcommand{\boldsymbol}[1]{\mbox{\boldmath $#1$}}

% Because html converters don't know tabularnewline  ====================
\providecommand{\tabularnewline}{\\}
\makeatother

% Hyper References  =====================================================
% Cria links entre referências e citações e as respectivas fórmulas, 
% figuras, tabelas, tópicos da bibliografia, seções, etc.
\usepackage{hyperref}
\hypersetup{
    colorlinks=true,
    linkcolor=blue,
    filecolor=magenta,      
    urlcolor=cyan,
%    pdftitle={Overleaf Example},
    pdfpagemode=FullScreen,
    }
%\hypersetup{
%citebordercolor=0 .75 1,
%linktocpage=true
%,backref=true
%}

%   TRADUÇÃO DO REVTEX4  %%%%%%%%%%%%%%%%%%%%%%%%%%%%%%%%%%%%%%%%%%%%%%%%%%

% Tradução para o Português de palavras geradas automaticamente  =======
\makeatletter
\def\tocname{Sumário}
\def\andname{e}
\def\Dated@name{Data: }
\def\lofname{Lista de Figuras}
\def\lotname{Lista de Tabelas}
\def\refname{Referências}
\renewcommand*\email[1][E-mail: ]{\begingroup\sanitize@url\@email{#1}}%
\def\today{\number\day \space de \space \ifcase\month\or janeiro\or fevereiro\or
mar{\c c}o\or abril\or maio\or junho\or julho\or agosto\or setembro\or outubro\or 
novembro\or dezembro\fi \space de \space \number\year}
\makeatother


%  ESQUELETO  %%%%%%%%%%%%%%%%%%%%%%%%%%%%%%%%%%%%%%%%%%%%%%%%%%%%%%%%%%%%%
\begin{document}

% Título do Relatório  ================================================
\title{Modelo \LaTeX \, para confecção de um relatório do IYPT-Brasil: uma apresentação bela e profissional}

% Nomes e E-mails dos Autores  ========================================
\author{Fabricio Marques do Carmo}
\email{fabriciomarques@if.usp.br}

% Replique o código acima para inserir mais autores!


% Nome da Equipe e Instituição  =======================================
\affiliation{Equipe: The Organizers, Sociedade Brasileira de Física}

% Coloque o nome da sua equipe e da sua escola no campo acima!


% Data  ===============================================================
\date{\today}  % O comando ``\today'' insere automaticamente a data do
               % momento em que o relatório é compilado.
               %
               % Para usar uma data diferente da atual, basta inserir
               % manualmente.

% Folha de Rosto  =====================================================
\begin{abstract}{
    \input{00-resumo.tex}  % Resumo
  }\end{abstract}

\maketitle    % Determina que a página do resumo é a página que contém o
              % título e insere título, nomes dos autores, data, etc.

\setstretch{2} % Seleciona ``doublespacing'' como espaçamento padrão 
               % entre linhas conforme regulamento IYPT-Brasil.
               % Mesmo quando o LaTeX acusa erro, geralmente funciona.
               
\newpage      % Inicia uma nova página de modo que título, nomes de 
              % autores e resumo ficam na primeira página e o relatório 
              % propriamente dito começa na página seguinte.

% Seções  =============================================================
\input{01-introducao.tex}   % Introdução
\input{02-fundamentos.tex}  % Fundamentos Teóricos
\input{03-descricao.tex}    % Descrição Experimental
\input{04-analise.tex}      % Analise dos Dados
\input{05-conclusao.tex}    % Conclusão
\input{06-bibliografia.tex} % Bibliografia
%\newpage                    % para iniciar os apêndices uma página depois
\input{07-apendices.tex}    % Apêndices

\end{document}

%%%%%%%%%%%%%%%%%%%%%%%%%%%%%%%%%%%%%%%%%%%%%%%%%%%%%%%%%%%%%%%%%%%%%%%%%

\end{document}